Table~\ref{tab:p4-errors-full} extends the error distribution to all ten runs.

\begin{table}[H]
\centering
\small
\begin{tabular}{lrrrrrr}
\toprule
\textbf{Run} & \textbf{No Ans.} & \textbf{Fmt Only} & \textbf{No Reas.} & \textbf{Arith.} & \textbf{Wrong Set.} & \textbf{Gibb.} \\
\midrule
Baseline & 318 (32.1\%) & 382 (38.6\%) & 0 & 108 (10.9\%) & 179 (18.1\%) & 3 (0.3\%) \\
\midrule
+ A & 163 (18.5\%) & 153 (17.4\%) & 0 & 168 (19.1\%) & 390 (44.3\%) & 6 (0.7\%) \\
+ B & 278 (28.7\%) & 302 (31.1\%) & 0 & 137 (14.1\%) & 249 (25.7\%) & 3 (0.3\%) \\
A+B & 205 (22.3\%) & 183 (19.9\%) & 0 & 178 (19.3\%) & 351 (38.1\%) & 3 (0.3\%) \\
\midrule
+ C & 89 (7.7\%) & 97 (8.4\%) & 0 & 216 (18.8\%) & 437 (38.0\%) & 2 (0.2\%) \\
+ D & 211 (18.1\%) & 237 (20.4\%) & 0 & 158 (13.6\%) & 311 (26.7\%) & 1 (0.1\%) \\
+ E & 63 (5.5\%) & 76 (6.7\%) & 0 & 239 (21.0\%) & 448 (39.4\%) & 2 (0.2\%) \\
C+D & 87 (7.6\%) & 74 (6.4\%) & 0 & 204 (17.7\%) & 453 (39.4\%) & 3 (0.3\%) \\
C+D+E & 41 (3.6\%) & 49 (4.3\%) & 0 & 220 (19.2\%) & 493 (43.0\%) & 0 \\
\midrule
Comb.\ All & 89 (7.9\%) & 54 (4.8\%) & 0 & 239 (21.2\%) & 466 (41.3\%) & 2 (0.2\%) \\
\bottomrule
\end{tabular}
\caption{Mistake type distribution across all ten runs. Percentages are within each run's failed set (problems where all 8 samples failed).}
\label{tab:p4-errors-full}
\end{table}

Number grounding (E) produces the most extreme format-failure reduction of any individual reward: ``no answer'' drops to 5.5\% ($-$26.6pp from baseline) and ``format only'' to 6.7\% ($-$31.9pp).
This surpasses even coherence (C), suggesting that number grounding's signal---rewarding responses that use the question's numbers---implicitly encourages structured mathematical responses that include parseable answers.

The C+D+E combination achieves the lowest format-failure rates of any configuration: 3.6\% ``no answer'' and 4.3\% ``format only'', with zero gibberish.
Its dominant failure mode is ``wrong setup'' at 43.0\%, the highest of any configuration, indicating that the grounding reward stack nearly eliminates surface-level failures, leaving only fundamental reasoning errors.
